\documentclass[FGBook.tex]{subfiles}

\begin{document}\changefontsize{8pt}

\proza{Šifra Mistra Leonarda / Da Vinciho Kód}{DAN BROWN}[TODO]

\noindent\begin{wrapfigure}{r}{0.25\textwidth}
\tiny

\subwection{Ukázka z díla:}
\setlength{\parindent}{3pt}
Langdon nemohl odtrhnout oči od fotografie, kterou držel v ruce, a pocítil neurčitý strach.
Obrázek byl příšerný a hluboce podivný, vyvolával zneklidňující dojem jakéhosi \textit{déjá-vu}.
Více než před rokem Langdon obdržel fotografii těla s~podobnou žádostí o~pomoc.
O čtyřiadvacet hodin později ve Vatikánu málem přišel o život.
Tento snímek byl naprosto odlišný, a přece na něm bylo něco povědomého.
Policista se podíval na hodinky.
\enquote{Můj kapitán vás očekává, pane.}
Langdon ho sotva slyšel.
Očima se vpíjel do obrázku.
\enquote{Ten symbol tady, a~to, jak~podivně je~to~tělo…} \texthl{\enquote{Naaranžováno?}} pomohl mu policista.
Langdon přikývl a~projel jím~chlad.
\enquote{Neumím~si~představit, kdo může tohle udělat jinému člověku.}
Policista vypadal zachmuřeně.
\enquote{Vy to nechápete, pane Langdone. To,~co~vidíte na té fotografii…} Odmlčel se.
\textithl{\enquote{Monsieur~Sauniére~si~to~udělal~sám.}}
\end{wrapfigure}

\wection{LITERÁRNÍ TEORIE}

\subwection{Literární druh a žánr:}
\noindent Prozaické epické dílo, \texthl{dobrodružný román},
\texthl{konspirační thriller} s detektivními prvky.

\subwection{Literární směr:}
\noindent Kniha patří do populární (až konzumní) literatury
\texthl{postmodernismu}, která kombinuje fakta a fikci.
Někdy se řadí jako samostatný směr navazující na postmodernismus.

\subwection{Jazyk a slovní zásoba:}
\noindent Dynamický a čtivý text s důrazem na akčnost.
Text je naplněn odbornými termíny,
převážně z okruhů \texthl{umění}, \texthl{historie} a \texthl{náboženství},
ale ve srozumitelných kontextech.)
Z cizích jazyků se zde vyskytuje \texthl{francouzština} a~\texthl{latina}.
Teabing zase hovoří staroangličtinou.

% \subwection{Figury:}

% \subwection{Tropy:}

\subwection{Stylistická charakteristika textu:}
\noindent Autor využívá spíše kratších vět, pro udržení dynamiky.
Často se vyskytují zvraty.

\subwection{Postavy:}
\noindent 
\textschl{Robert Langdon --} profesor náboženského symbolismu.
Je znenadání zavolán na místo činu, kdy se dovídá, že zemřel jeho kolega z oboru.
Stejně jako se mu to \texthl{již stalo v minulosti}, je zatažen do víru událostí,
se kterým má pramálo společného.
Či v to alespoň věří.
S tím, jak odkrývá záhady kolem Jacquese Saunièra, tak se dozvídá i důvod,
pro jeho zapojení do událostí.\\
\textschl{Sophie Neveuová --} policejní kryptoložka,
se kterou se Langdon setkává na místě činu.
Varuje ho i před skutečností, že na něj Fache chce vraždu hodit.
Přiznává, že je vnučkou pana Saunièra a že má klíč k vyřešení hádanek.
Ani ona ale netuší rozsah konspirace, kterou se chystají odhalit.\\
\textschl{Sir Leigh Teabing / Učitel --} z prvu se tváří jako pomocník hlavní dvojice.
Bohatý britský historik, gentleman ze staré školy,
který je nadšen z~představy hledání Svatého Grálu a nezištně jim pomáhá.
Později se ukazuje, že to on, v roli takzvaného Učitele, dal věci do pochybu.
Neplánoval tedy, že ho Langdon sám vyhledá, tím mu to značně usnadnil,
ale to on může za vraždy a pro posedlost Grálem se obrátí i proti Langdonovi.\\
\textschl{Silas --} tragická postava.
Retrospektivně sledujeme jeho mládí.
Jak ho týrali.
Jak se ho ujal Aringarosa.
Jak zasvětil celý svůj život službě muži,
který vůči němu jako jediný projevil lítost.
Aringarosa svého \enquote{Božího válečníka} svěřil do služby Učiteli.
A tak se Silas hned při několika příležitostech střetl s Langdonem.\\
\textschl{Bezu Fache --} kapitán francouzské policie.
Člen Opus Dei.
Biskup Aringarosa ho přesvědčil o Langdonově vině
a svěřil mu za úkol jeho dopadení.
Jinak poctivý policista, na Langdonovu stopu se nedal kvůli korupci,
ale kvůli důveře ve slovo Božího služebníka.\\
\textschl{Jacques Saunière --} kurátor Louvru,
velmistr \texthl{Převorství sionského},
dědeček Sophie Neveuové,
historik a symbolog.
Se svojí vnučkou měl složitý vztah.
Langdona pověřil zachováním tajemství poté, co věřil,
že tajemství, aspoň jeho část, už odhalil, díky poznámce v jeho knize,
kterou poslal Saunièrovi na přečtení.\\
\textschl{Biskup Manuel Aringarosa --} biskup Opus Dei.
Pro svou organizaci je schopen zajít do extrémů.
Je manipulativní (Bezu Fache),
stejně se ale \texthl{nechá zmanipulovat} (Učitelem).

\subwection{Děj:}
\noindent Ocitáme se v Louvru.

\subwection{Kompozice, prostor a čas:}
\noindent 
xXx

% \subwection{Význam sdělení:}

\wection{LITERÁRNÍ HISTORIE}

% \subwection{Politická situace:}

% \subwection{Základní principy fungování společnosti:}

% \subwection{Kontext dalších druhů umění:}

% \subwection{Kontext literárního vývoje:}

\subwection{Autor {\ssmall -- život autora:}}
\noindent 
xXx

\subwection{Vlivy na dané dílo:}
\noindent xXx

\subwection{Další autorova tvorba:}
\noindent 
Mezi nejznámnější díla patří \textsc{xXx}

% \subwection{Jak dílo inspirovalo další vývoj literatury:}

\wection{LITERÁRNÍ KRITIKA}

\subwection{Dobové vnímání díla a jeho proměny:}
\noindent xXx

\subwection{Aktuálnost tématu a zpracování tématu:}
\noindent xXx

\subwection{Kontrola faktů:}
\noindent
\textschl{Opus Dei --} xXx.\\
\textschl{Převorství Sionské --} xXx.

\wection{OSOBNÍ NÁZOR}
\noindent 
xXx

\vfill

\noindent\begin{minipage}{\textwidth}
    {\textcolor{\wpagecolor}{\rule{\linewidth}{0.4pt}}
    \footnotesize
    \textbfhl{dobrodružný román --} xXx;
    \textbfhl{konspirační thriller --} xXx;
    \textbfhl{populární literatura --} xXx;
    \textbfhl{konzumní literatura --} xXx 
    }
\end{minipage}

\newpage

\end{document}